% Part: first-order-logic
% Chapter: introduction
% Section: sentences

\documentclass[../../../include/open-logic-section]{subfiles}

\begin{document}

\olfileid{fol}{int}{snt}

\olsection{\usetoken{P}{sentence}}

Saiki kita wis nduweni (rancangan) definisi relasi nyukupi (``bener
ing'') kanggo !!{structure}s lan !!{formula}s. Nanging, definisi iki
mbutuhake perangan tambahan---yaiku pangaji !!a{variable}---dene kang
kita karepake yaiku definisi !!{sentence}s. Kepriye carane nyingkirake
pangaji, lan apa iku !!{sentence}s?

Sampeyan mbokmenawa isih kelingan rembugan ing pambuka pisanan marang
logika formal ngenani relasi antarane !!{variable}s lan kuantor. Kuantor
mesthi diterusake dening !!a{variable}, banjur ing perangan
!!{sentence} kang dikenani kuantor kasebut (``cakupane''), kita
mangerteni yen !!{variable} kasebut ``kaiket'' dening kuantor mau. Ing
!!{formula}s, saben !!{variable} ora kudu nduweni kuantor kang cocog,
lan !!{variable}s tanpa kuantor kang cocog iku ``bebas'' utawa ``ora
kaiket.'' Kita bakal nganggep !!{sentence}s minangka kabeh !!{formula}s
kang ora nduweni !!{variable}s bebas.

Maneh, gagasan intuisi ngenani kapan muncule !!a{variable} ing
!!a{formula}~$!A$ iku kaiket, kuantor endi kang ngiket, lan kapan
muncule bebas, ora angel dimangerteni. Sampeyan bisa uga wis sinau cara
kanggo mriksa iki, mbokmenawa kanthi ngetung kurung. Kita kudu tetep
nyuwun definisi kang trep---lan amarga wis nemtokake !!{formula}s kanthi
induksi, kita uga bisa menehi definisi kanthi induksi tumrap saben
muncule !!a{variable}~$x$ kang bebas lan kaiket ing
!!a{formula}~$!A$. Tuladhane, kanggo basa prasaja kita, definisine bisa
kaya mangkene:
\begin{enumerate}
  \item Yen $!A$ iku atom, saben muncule $x$ ing formula kasebut bebas
  (yaiku, muncule $x$ ing~$\Atom{\Obj P}{x}$ iku bebas).
  \item Yen $!A$ awujud $\lnot !B$, banjur sawijining muncule~$x$
  ing~$\lnot !B$ bebas yen lan mung yen muncule~$x$ kang cocog bebas
  ing~$!B$ (yaiku, saben muncule variabel bebas ing~$
  !B$ persis padha karo muncule kang cocog ing~$\lnot !B$).
  \item Yen $!A$ awujud $(!B \land !C)$, banjur sawijining muncule~$x$
  ing~$(!B \land !C)$ bebas yen lan mung yen muncule~$x$ kang cocog
  bebas ing~$!B$ utawa ing~$!C$.
  \item Yen $!A$ awujud $\lexists[x][!B]$, banjur ora ana muncule
  $x$ ing~$!A$ kang bebas; yen formula kasebut awujud $\lexists[y][!B]$ kanthi
  $y$ minangka !!{variable} kang beda saka~$x$, banjur sawijining
  muncule~$x$ ing $\lexists[y][!B]$ bebas yen lan mung yen muncule~$x$
  kang cocog bebas ing~$!B$.
\end{enumerate}

Sawise nduweni definisi kang trep tumrap muncule variabel bebas lan
kaiket, kita mung perlu kandha: !!a{sentence} yaiku sembarang
!!{formula} tanpa muncule !!{variable}s bebas.

\end{document}
